\documentclass[../../../include/open-logic-section]{subfiles}

\begin{document}

\olfileid[hi]{sth}{cardinals}{cardsasords}
\olsection{क्रमसूचकों के रूप में कार्डिनल}

वास्तव में कार्डिनलों का हमारा सिद्धांत क्रमसूचकों के हमारे सिद्धांत का
(निर्लज्जतापूर्वक) उपयोग करेगा। अर्थात् हम कार्डिनलों को कुछ विशिष्ट
क्रमसूचकों के रूप में ही परिभाषित करेंगे। विशेषतः हम निम्न परिभाषा देंगे:

\begin{defn}\ollabel{defcardinalasordinal}
यदि $A$ को सुव्यवस्थित किया जा सकता है, तो $\card{A}$ वह लघुतम क्रमसूचक
$\gamma$ है जिसके लिए $\cardeq{A}{\gamma}$। किसी भी क्रमसूचक $\gamma$ के
लिए हम कहते हैं कि $\gamma$ एक \emph{कार्डिनल} है तभी जब
$\gamma = \card{\gamma}$।
\end{defn}

हमने अभी ``$A$ को सुव्यवस्थित किया जा सकता है'' वाक्यांश का उपयोग किया।
गणित में लगभग हमेशा की तरह यहाँ यह मोडल कथन किसी अस्तित्वपरक दावे का केवल
अनौपचारिक रूप है: ``$A$ को सुव्यवस्थित किया जा सकता है'' कहने का अर्थ केवल
इतना है कि ``कोई ऐसा संबंध है जो $A$ को सुव्यवस्थित करता है''।

किंतु \olref{defcardinalasordinal} में एक अड़चन है। हम चाहेंगे कि
\emph{प्रत्येक} समुच्चय का कोई आकार हो, अर्थात् प्रत्येक $A$ के लिए
$\card{A}$ का अस्तित्व हो। किंतु हमारी अभी दी गई परिभाषा एक सशर्त वाक्य से
आरंभ होती है: ``\emph{यदि} $A$ को सुव्यवस्थित किया जा सकता है\ldots''। यदि
कोई ऐसा समुच्चय $A$ हो जिसे सुव्यवस्थित नहीं किया जा सकता, तो हमारी परिभाषा
$\card{A}$ नामक किसी वस्तु को परिभाषित ही नहीं करेगी।

अतः \olref{defcardinalasordinal} का उपयोग करने के लिए हमें यह प्रत्याभूति
चाहिए कि प्रत्येक समुच्चय को सुव्यवस्थित किया जा सकता है। दुर्भाग्य से
$\ZF$ में यह प्रत्याभूति उपलब्ध नहीं है। इसलिए यदि हम
\olref{defcardinalasordinal} का उपयोग करना चाहते हैं, तो कोई नया अभिगृहीत
जोड़ने के अतिरिक्त कोई विकल्प नहीं है, जैसे:
\begin{axiom}[Well-Ordering]
प्रत्येक समुच्चय को सुव्यवस्थित किया जा सकता है।
\end{axiom}
हम \olref[choice][]{chap} में चर्चा करेंगे कि सुव्यवस्था अभिगृहीत स्वीकार्य
है या नहीं। किंतु अब से हम इसका मुक्त रूप से उपयोग करेंगे। इसके उपयोग से यह
सिद्ध करना बहुत सीधा है कि (\olref{defcardinalasordinal} में परिभाषित)
कार्डिनलों का अस्तित्व है और वे सुव्यवस्थित ढंग से व्यवहार करते हैं:

\begin{lem}\ollabel{lem:CardinalsExist}
प्रत्येक समुच्चय $A$ के लिए:
\begin{enumerate}
	\item\ollabel{cardaexists} $\card{A}$ का अस्तित्व है और वह अद्वितीय है;
	\item\ollabel{cardaapprox}  $\cardeq{\card{A}}{A}$;
	\item\ollabel{cardaidem}  $\card{A}$ एक कार्डिनल है, अर्थात्
	$\card{A} = \card{\card{A}}$;
\end{enumerate}
\end{lem}

\begin{proof}
$A$ को नियत करें। सुव्यवस्था से कोई सुव्यवस्था $\tuple{A, R}$ है।
\olref[ordinals][ordtype]{thmOrdinalRepresentation} से $\tuple{A,R}$
किसी अद्वितीय क्रमसूचक $\beta$ के समाकृतिक है। अतः
$\cardeq{A}{\beta}$। परिमितेतर आगमन से कोई अद्वितीय लघुतम क्रमसूचक
$\gamma$ है जिसके लिए $\cardeq{A}{\gamma}$। अतः $\card{A}=\gamma$,
जिससे \olref{cardaexists} और \olref{cardaapprox} स्थापित होते हैं।
\olref{cardaidem} स्थापित करने के लिए ध्यान दें कि यदि $\delta \in
\gamma$, तो $\gamma$ के हमारे चयन से $\cardless{\delta}{A}$, और इसलिए
$\cardless{\delta}{\gamma}$ भी, क्योंकि समसंख्यता एक तुल्यता संबंध है
(\olref[sfr][siz][equ]{equinumerosityisequi})। अतः
$\gamma=\card{\gamma}$।
%So, for reductio, suppose that there is some ordinal $\delta \in \gamma$ such that $\cardeq{\gamma}{\delta}$. Then, $\cardeq{\cardeq{A}{\gamma}}{\delta}$ so that $\cardeq{A}{\delta}$ by \olref[sfr][set][equ]{equinumerosityisequi}, which contradicts the choice of $\gamma$ as the least ordinal such that $\cardeq{A}{\gamma}$. 
\end{proof}

अगला परिणाम कैंटर के सिद्धांत की, और उसके अतिरिक्त भी, प्रत्याभूति देता है।
(ध्यान दें कि कार्डिनल अपना क्रम क्रमसूचकों से प्राप्त करते हैं, अर्थात्
$\cardfont{a} < \cardfont{b}$ तभी जब $\cardfont{a} \in \cardfont{b}$। इसे
सूत्रबद्ध करते समय जिन वस्तुओं को हम कार्डिनल जानते हैं, उनके लिए फ़्राक्टूर
अक्षरों का उपयोग करेंगे। यह पर्याप्त मानक प्रथा है। एक सामान्य विकल्प
यूनानी अक्षरों का उपयोग करना है, क्योंकि कार्डिनल क्रमसूचक होते हैं, किंतु
उन्हें वर्णमाला के मध्य से चुनना, जैसे $\kappa, \lambda$।):
\begin{lem}\ollabel{lem:CardinalsBehaveRight}
किन्हीं समुच्चयों $A$ और $B$ के लिए:
\begin{align*}
	\cardeq{A}{B} &\text{ तभी जब } \card{A} = \card{B}\\
	\cardle{A}{B} &\text{ तभी जब } \card{A} \leq \card{B}\\
	\cardless{A}{B}&\text{ तभी जब } \card{A} < \card{B}
\end{align*}
\end{lem}

\begin{proof}
हम दूसरे दावे की बाएँ-से-दाएँ दिशा सिद्ध करेंगे (अन्य स्थितियाँ समान हैं
और अभ्यास के लिए छोड़ी गई हैं)। अतः निम्न आरेख पर विचार करें:
\begin{center}
	\begin{tikzpicture}
	\node (nodea) {$A$};
	\node[right = 6em of nodea] (nodeb) {$B$};
	\node[below = 2em of nodea] (nodecarda) {$\card{A}$};
	\node[below = 2em of nodeb] (nodecardb) {$\card{B}$};
	\draw[->] (nodea)--(nodeb);
	\draw[<->] (nodea)--(nodecarda);
	\draw[<->] (nodeb)--(nodecardb);
	\draw[->, dashed] (nodecarda)--(nodecardb);
\end{tikzpicture}
\end{center}
दो सिरों वाले तीर !!{bijection}ों को दर्शाते हैं, जिनके अस्तित्व की
प्रत्याभूति \olref{lem:CardinalsExist} से मिलती है। $\cardle{A}{B}$ मानने
पर कोई !!a{injection} $A\to B$ है। अब $\card{A}$ से $A$, फिर $B$, और फिर
$\card{B}$ तक तीरों का अनुसरण करके हमें कोई !!a{injection}
$\card{A}\to\card{B}$ (टूटा हुआ तीर) मिलता है।
\end{proof}\noindent
\olref{lem:CardinalsBehaveRight} का उपयोग श्रेडर--बर्नस्टीन को फिर से
सिद्ध करने में भी कर सकते हैं। यह दावा है कि यदि $\cardle{A}{B}$ और
$\cardle{B}{A}$, तो $\cardeq{A}{B}$। हमने इसे
\olref[sfr][siz][sb]{thm:schroder-bernstein} के रूप में कहा था, किंतु पहले
इसे---कुछ प्रयास के साथ---\olref[sfr][infinite][card-sb]{sec} में सिद्ध
किया था। अब विचार करें:

\begin{proof}[श्रेडर--बर्नस्टीन का पुनः प्रमाण]
यदि $\cardle{A}{B}$ और $\cardle{B}{A}$, तो
\olref{lem:CardinalsBehaveRight} से $\card{A}\leq\card{B}$ और
$\card{B}\leq\card{A}$। अतः त्रिविभाजन और
\olref{lem:CardinalsBehaveRight} से $\card{A}=\card{B}$ और
$\cardeq{A}{B}$।
\end{proof}
\noindent
यद्यपि यह बहुत सरल प्रमाण है, यह अप्रत्यक्ष रूप से प्रतिस्थापन
(\olref[ordinals][ordtype]{thmOrdinalRepresentation} पाने के लिए) और
सुव्यवस्था (\olref{lem:CardinalsBehaveRight} की प्रत्याभूति के लिए) दोनों
पर निर्भर है। इसके विपरीत \olref[sfr][infinite][card-sb]{sec} का प्रमाण
कहीं अधिक स्वावलंबी था (वास्तव में उसे $\Zminus$ में पूरा किया जा सकता है)।

\end{document}
